\documentclass[12pt, a4paper]{article}

% --- Cấu hình Tiếng Việt và Font ---
\usepackage[utf8]{inputenc}
\usepackage[vietnamese]{babel}
\usepackage[T1]{fontenc}
\usepackage{lmodern}

% --- Các gói bổ trợ ---
\usepackage{amsmath, amssymb, amsthm}
\usepackage{graphicx}
\usepackage{hyperref}
\usepackage{xcolor}
\usepackage{geometry}
\usepackage{listings}
\usepackage{enumitem}
\usepackage{indentfirst} % Thụt đầu dòng đoạn văn đầu tiên

% --- Cấu hình lề trang ---
\geometry{
    a4paper,
    left=25mm,
    right=25mm,
    top=25mm,
    bottom=25mm
}

% --- Cấu hình hiển thị Code (JavaScript) ---
\definecolor{codegreen}{rgb}{0,0.6,0}
\definecolor{codegray}{rgb}{0.5,0.5,0.5}
\definecolor{codepurple}{rgb}{0.58,0,0.82}
\definecolor{backcolour}{rgb}{0.98,0.98,0.95}

\lstset{
    language=JavaScript,
    backgroundcolor=\color{backcolour},   
    commentstyle=\color{codegreen},
    keywordstyle=\color{blue},
    numberstyle=\tiny\color{codegray},
    stringstyle=\color{codepurple},
    basicstyle=\ttfamily\small,
    breakatwhitespace=false,         
    breaklines=true,                 
    captionpos=b,                    
    keepspaces=true,                 
    numbers=left,                    
    numbersep=5pt,                  
    showspaces=false,                
    showstringspaces=false,
    showtabs=false,                  
    tabsize=2,
    frame=lines
}

% --- Thông tin tiêu đề ---
\title{\textbf{KIẾN TRÚC HƯỚNG DỊCH VỤ} \\ \vspace{0.5cm} \Large Buổi 8: API Testing và Quality Assurance (QA)}

\date{}

\begin{document}

\maketitle

\section{Giới thiệu chung}
Trong Kiến trúc hướng dịch vụ (SOA), API đóng vai trò là "hợp đồng" giao tiếp giữa các thành phần. Việc đảm bảo chất lượng (Quality Assurance) cho API không chỉ dừng lại ở việc kiểm tra tính đúng đắn của dữ liệu mà còn phải đảm bảo tính sẵn sàng, hiệu năng và khả năng chịu tải của toàn bộ hệ thống.

\section{Kiến thức trọng tâm}

\subsection{Phân loại các loại hình Kiểm thử (Testing Types)}
Việc kiểm thử được phân tầng theo mô hình Kim tự tháp thử nghiệm (Testing Pyramid):

\begin{enumerate}[leftmargin=*]
    \item \textbf{Unit Test (Kiểm thử mức đơn vị):} 
    \begin{itemize}
        \item Tập trung vào các hàm (functions) hoặc phương thức đơn lẻ trong mã nguồn.
        \item Sử dụng kỹ thuật \textit{Mocking} để giả lập kết quả từ database hoặc các service khác, giúp test chạy độc lập và cực nhanh.
        \item Mục tiêu: Phát hiện lỗi logic ngay tại thời điểm viết code.
    \end{itemize}

    \item \textbf{Integration Test (Kiểm thử tích hợp):} 
    \begin{itemize}
        \item Kiểm tra sự tương tác giữa các module đã được kết nối (ví dụ: API kết nối với Database thật).
        \item Đảm bảo luồng dữ liệu (data flow) không bị sai lệch khi đi qua các lớp kiến trúc (Controller $\rightarrow$ Service $\rightarrow$ Repository).
    \end{itemize}

    \item \textbf{Performance Test (Kiểm thử hiệu năng):} 
    \begin{itemize}
        \item \textit{Load Testing:} Giả lập số lượng người dùng truy cập đồng thời ở mức dự kiến để đo phản ứng của hệ thống.
        \item \textit{Stress Testing:} Đẩy lượng truy cập vượt quá giới hạn thiết kế để xác định điểm gãy (breaking point) của hệ thống.
        \item \textit{Spike Testing:} Kiểm tra khả năng xử lý khi lượng truy cập tăng đột biến trong thời gian ngắn.
    \end{itemize}
\end{enumerate}

\subsection{Công cụ hỗ trợ chuyên dụng}
\begin{itemize}
    \item \textbf{Postman:} Môi trường phát triển API toàn diện. Cho phép tạo các Collection, thiết lập biến môi trường (Environment Variables) và viết script kiểm thử tự động.
    \item \textbf{Newman (CLI for Postman):} Cho phép thực thi các bộ test Postman trực tiếp trên terminal. Đây là công cụ quan trọng để tích hợp vào luồng tự động hóa CI/CD (như Jenkins hoặc GitHub Actions).
    \item \textbf{k6 \& JMeter:} Các công cụ chuyên dụng cho Performance Testing. Trong đó \textbf{k6} được ưa chuộng nhờ việc viết kịch bản bằng JavaScript và khả năng mở rộng tốt trong môi trường cloud.
\end{itemize}

\section{Kỹ năng thực hành chi tiết}

\subsection{Xây dựng bộ Test tự động (Automated API Testing)}
Kịch bản test tự động trong Postman được viết trong tab \textit{Tests} bằng ngôn ngữ JavaScript. Dưới đây là bộ test đầy đủ cho một endpoint:

\begin{lstlisting}[caption=Kịch bản kiểm thử tự động toàn diện]
// 1. Kiểm tra trạng thái phản hồi (Status Code)
pm.test("Xác nhận mã trạng thái là 200 OK", function () {
    pm.response.to.have.status(200);
});

// 2. Kiểm tra thời gian phản hồi (Response Time)
pm.test("Thời gian phản hồi phải dưới 300ms", function () {
    pm.expect(pm.response.responseTime).to.be.below(300);
});

// 3. Kiểm tra Header của phản hồi
pm.test("Content-Type phải là application/json", function () {
    pm.response.to.have.header("Content-Type", "application/json; charset=utf-8");
});

// 4. Kiểm tra cấu trúc và nội dung dữ liệu JSON
pm.test("Dữ liệu trả về đúng định dạng", function () {
    var jsonData = pm.response.json();
    
    // Kiểm tra các trường bắt buộc
    pm.expect(jsonData).to.have.property('id');
    pm.expect(jsonData.id).to.be.a('number');
    
    // Kiểm tra tính chính xác của dữ liệu
    pm.expect(jsonData.username).to.eql("admin");
});
\end{lstlisting}

\subsection{Đo lường và phân tích hiệu năng (Performance Measurement)}
Khi thực hiện đo hiệu năng, sinh viên cần tập trung vào các chỉ số kỹ thuật sau:

\subsubsection{Các chỉ số về Thời gian (Response Time Metrics)}
\begin{itemize}
    \item \textbf{Average (Trung bình):} Tổng thời gian chia cho số request.
    \item \textbf{p95, p99 (Percentiles):} Ý nghĩa cực kỳ quan trọng. Ví dụ p95 = 500ms nghĩa là 95\% người dùng có thời gian phản hồi dưới 500ms. Chỉ có 5\% bị chậm hơn mức này. p99 phản ánh trải nghiệm của nhóm người dùng gặp tệ nhất.
\end{itemize}

\subsubsection{Các chỉ số về Độ ổn định (Stability Metrics)}
\begin{itemize}
    \item \textbf{Error Rate (Tỷ lệ lỗi):} Số lỗi (thường là mã 5xx) trên tổng số request gửi đi. Một hệ thống tốt phải có tỷ lệ lỗi xấp xỉ 0\% dưới tải trọng dự kiến.
    \item \textbf{Throughput (Thông lượng):} Số lượng request mà API có thể xử lý thành công trong 1 giây (RPS - Requests Per Second).
\end{itemize}

\section{Kết luận}
Việc thực hiện tốt các kỹ năng trong Buổi 8 giúp đảm bảo hệ thống SOA hoạt động ổn định, tin cậy. Kết hợp giữa việc kiểm thử tự động với Postman/Newman và đo hiệu năng với k6 giúp giảm thiểu rủi ro khi triển khai dịch vụ thực tế.

\end{document}